\chapter{Nuclear Fusion}

\section{Nucleosynthesis}
 
\subsection{Processes inside of a star}
 At the core of stars, lighter elements get converted into heavier elements through nuclear fusion processes. For instance, in newly born star, Hydrogen get converted into Helium which produces immense amount of thermal energy which illuminates the whole surface of the collapsing dust cloud. For millions to billions years, the star passes through different evolutionary triads, primarily influenced by its mass.  
 
 \vspace{1 cm}
 
 Proton-Proton Reactions: Efficient for 1 $M_0$ star's core. \\
 
 \hspace{3cm} Branch I: 84\% Helium production, at 10-18 MK
 
 \begin{center}
 	\ce{H ->[p][e+,$\nu_e$]D ->[H][ $\gamma$ ]  He ->[He][2H] He } 	
 \end{center}
 \vspace{-2mm}
 
 \hspace{5.9cm} {\scriptsize 1 \hspace{0.8cm} 2 \hspace{0.9cm} 3 \hspace{0.9cm} 4   \hfill (atomic mass) \\
 	\vspace{-3mm}
 	\hspace{6.7cm} 1442 \hspace{0.5cm} 5493 \hspace{0.5cm} 12859 \hspace{0.9cm} \hfill (Energy release in KeV)}
 
 
 
 \vspace{5mm}
 \hspace{3cm} Branch II: 16.68 \% Helium production, at 18-25 MK
 
 \begin{center}
 	\ce{He ->[He][$\gamma$]Be ->[e-][ $\nu_e$]  Li ->[H] He }	
 \end{center}
 \vspace{-2mm}
 \hspace{6.5cm} {\scriptsize 4 \hspace{0.9cm} 7 \hspace{0.9cm} 7 \hspace{1.1cm} 4  \hfill (atomic mass) \\
 	\vspace{-3mm}
 	\hspace{6.8cm} 1590 \hspace{0.3cm} 861/383 \hspace{0.3cm} 17350 \hspace{0.9cm} \hfill (Energy release in KeV)}
 
 
 \vspace{4cm}
 
 \hspace{3cm} Branch III: 0.02\% Helium production, above 25 MK
 
 \begin{center}
 	\ce{He ->[He][$\gamma$]Be ->[H][ $\gamma$]  B ->[][e+,$\nu$_e] Be -> 2He }	
 \end{center}
 \vspace{-2mm}
 \hspace{5.6cm} {\scriptsize 4 \hspace{0.8cm} 7 \hspace{0.9cm} 8 \hspace{1.0cm} 8 \hspace{1.0cm} 4   \hfill (atomic mass) \\
 	\vspace{-2mm}
 	
 	\hspace{5.5cm} 1590 \hspace{0.3cm}  \hfill (Energy release in KeV)}
 
 \vspace{1cm}
 
 
 \hspace{3cm} Branch IV: 0.000002\% Helium generation in Sun 
 
 \begin{center}
 	\ce{He ->[H][e+,$\nu$_e] He}	
 \end{center}
 \vspace{-2mm}
 \hspace{7.8cm} {\scriptsize 3 \hspace{0.9cm} 4   \hfill (atomic mass) \\
 	\vspace{-2mm}
 	
 	\hspace{7.6cm} 19795 \hspace{0.3cm}  \hfill (Energy release in KeV)}
 
 \vspace{1cm}
 
 
 Triple-alpha Process:\\
 
 \begin{center}
 	\ce{He ->[He] Be ->[He][2$\gamma$] C }
 \end{center}
 
 \vspace{-2mm}
 
 \hspace{6.6cm} {\scriptsize 4 \hspace{0.8cm} 8 \hspace{0.8cm} 12 \hfill (atomic mass) \\
 	\vspace{-2mm}
 	
 	\hspace{6.9cm} -91.8 \hspace{0.4cm}  7367 \hfill (Energy release in KeV)}
 
 \vspace{1cm}
 
 
 \vspace{1cm}
 
 CNO Cycle: Catalytic cycle dominant in 1.3 $M_0$ + stars\\
 
 \hspace{3 cm} Branch I: \\	
 \begin{center}
 	\ce{C ->[H][$\gamma$] N ->[][e+, $\nu$]C ->[H][$\gamma$] N ->[H][$\gamma$] O ->[][e+, $\nu$] N ->[H][He]C}
 \end{center}
 \vspace{-2mm}
 \hspace{4.7cm} {\scriptsize 12 \hspace{0.6cm} 13 \hspace{0.9cm} 13 \hspace{0.6cm} 14 \hspace{0.7cm} 15 \hspace{0.9cm} 15 \hspace{0.7cm} 12 \hfill (atomic mass) \\
 	\vspace{-2mm}
 	
 	\hspace{4.6cm} 1950 \hspace{0.4cm}  1200 \hspace{0.5cm}  7540 \hspace{0.4cm}  7350 \hspace{0.5cm}  1730 \hspace{0.5cm} 4960 \hfill (Energy release in KeV)}
 
 \vspace{1cm}
 
 
 \hspace{3 cm} Branch II: \\	
 \begin{center}
 	\ce{N ->[H][$\gamma$] O ->[H][$\gamma$]F ->[ ][e+, $\nu$_e] O ->[H][He] N ->[H][$\gamma$]O ->[][e+, \nu_e]N}
 \end{center}
 \vspace{-2mm}
 \hspace{4.5cm} {\scriptsize 15 \hspace{0.7cm} 16 \hspace{0.7cm} 17 \hspace{0.9cm} 17 \hspace{0.7cm} 14 \hspace{0.8cm} 15 \hspace{0.9cm} 15 \hfill (atomic mass) \\
 	\vspace{-2mm}
 	
 	\hspace{4.4cm} 12130 \hspace{0.4cm}  600 \hspace{0.5cm}  2760 \hspace{0.6cm}  1190 \hspace{0.5cm}  7350 \hspace{0.5cm} 2750 \hfill (Energy release in KeV)}
 
 \vspace{1cm}
 
 
 \hspace{3 cm} Branch III: \\	
 \begin{center}
 	\ce{O ->[H][$\gamma$] F ->[ ][e+, $\nu$_e]O ->[H][He] N ->[H][$\gamma$] O ->[H][$\gamma$]F ->[][e+, \nu_e]O}
 \end{center}
 \vspace{-2mm}
 \hspace{4.7cm} {\scriptsize 17 \hspace{0.6cm} 18 \hspace{0.7cm} 18 \hspace{0.5cm} 15 \hspace{0.6cm} 16 \hspace{0.8cm} 17 \hspace{0.5cm} 17 \hfill (atomic mass) \\
 	
 	\vspace{-2mm}
 	\hspace{4.6cm} 5610 \hspace{0.4cm}  1656 \hspace{0.3cm}  3980 \hspace{0.2cm}  12130 \hspace{0.3cm}  0.600 \hspace{0.3cm} 2760 \hfill (Energy release in KeV)}
 
 \vspace{1cm}
 
 
 \hspace{3 cm} Branch IV: \\	
 \begin{center}
 	\ce{O ->[H][$\gamma$] F ->[H][He]O ->[H][$\gamma$] F ->[ ][e+, $\nu$_e]O ->[H][$\gamma$] F ->[ ][e+,$\nu$_e]O}
 \end{center}
 \vspace{-2mm}
 \hspace{4.7cm} {\scriptsize 18 \hspace{0.6cm} 19 \hspace{0.6cm} 16 \hspace{0.7cm} 17 \hspace{0.9cm} 17 \hspace{0.7cm} 18 \hspace{1.0cm} 18 \hfill (atomic mass) \\
 	
 	\vspace{-2mm}
 	
 	\hspace{4.6cm} 7994 \hspace{0.4cm}  8114 \hspace{0.4cm}  600 \hspace{0.4cm}  2760 \hspace{0.7cm}  5610 \hspace{0.5cm} 1656 \hfill (Energy release in KeV)}
 
 \vspace{1cm}
 
 
 
 \hspace{3 cm} HCNO-I \\	
 \begin{center}
 	\ce{C ->[H][$\gamma$] N ->[H][$\gamma$]O ->[ ][e+, $\nu$_e]N ->[H][$\gamma$] O ->[ ][e+, $\nu$_e]N ->[H][He]C}
 \end{center}
 \vspace{-2mm}
 \hspace{4.6cm} {\scriptsize 12 \hspace{0.6cm} 13 \hspace{0.7cm} 13 \hspace{1.0cm} 14 \hspace{0.7cm} 15 \hspace{0.9cm} 16 \hspace{0.7cm} 12 \hfill (atomic mass) \\
 	\vspace{-2mm}
 	
 	\hspace{4.5cm} 1950 \hspace{0.4cm}  4630 \hspace{0.5cm}  5140 \hspace{0.6cm}  7350 \hspace{0.5cm}  2750 \hspace{0.6cm} 4960 \hfill (Energy release in KeV)}
 
 \vspace{1cm}
 
 
 
 \hspace{3 cm} HCNO-II \\	
 \begin{center}
 	\ce{N ->[H][$\gamma$] O ->[H][$\gamma$]F ->[H][$\gamma$]Ne ->[ ][e+,$\nu$_e]F ->[H][He] O ->[ ][e+,$\nu$_e]N}
 \end{center}
 \vspace{-2mm}
 \hspace{4.5cm} {\scriptsize 15 \hspace{0.7cm} 16 \hspace{0.7cm} 17 \hspace{0.7cm} 18 \hspace{1.1cm} 18 \hspace{0.6cm} 15 \hspace{0.9cm} 15 \hfill (atomic mass) \\
 	
 	\vspace{-2mm}
 	
 	\hspace{4.4cm} 12130 \hspace{0.5cm}  600 \hspace{0.4cm}  3920 \hspace{0.6cm}  4440 \hspace{0.6cm}  2880 \hspace{0.5cm} 2750 \hfill (Energy release in KeV)}
 
 \vspace{1cm}
 
 
 
 \hspace{3 cm} HCNO-III 	
 \begin{center}
 	\ce{F ->[H][$\gamma$] Ne ->[][e+,$\nu$_e]F ->[H][He]O ->[H][$\gamma$]F ->[H][$\gamma$]Ne ->[ ][e+,$\nu$_e]F}
 \end{center}
 \vspace{-2mm}
 \hspace{4.5cm} {\scriptsize 18 \hspace{0.7cm} 19 \hspace{1.0cm} 19 \hspace{0.7cm} 16 \hspace{0.7cm} 17 \hspace{0.7cm} 18 \hspace{1.0cm} 18 \hfill (atomic mass) \\
 	
 	\vspace{-2mm}
 	
 	\hspace{4.4cm} 6410 \hspace{0.7cm}  3320 \hspace{0.5cm}  8110 \hspace{0.5cm}  600 \hspace{0.5cm}  3920 \hspace{0.6cm} 4440 \hfill (Energy release in KeV)}
 
 \vspace{1cm}
 
 
 
Heavy Elements: \\

\hspace{3 cm} Alpha Ladder: \\


\begin{center}
	\ce{C ->[He][$\gamma$] O ->[He][$\gamma$] Ne  ->[He][$\gamma$] Mg  ->[He][$\gamma$] Si  ->[He][$\gamma$] S  ->[He][$\gamma$] Ar  ->[He][$\gamma$] Ca  ->[He][$\gamma$]  Ti  ->[He][$\gamma$]  Cr  ->[He][$\gamma$]  Fe  ->[He][$\gamma$]  Ni}
\end{center}
\vspace{-2mm}
\hspace{1.5cm}{\scriptsize 12 \hspace{0.7cm} 16 \hspace{0.7cm} 20 \hspace{0.8cm} 24 \hspace{0.8cm} 28 \hspace{0.8cm} 32 \hspace{0.8cm} 36 \hspace{0.8cm} 40 \hspace{0.8cm} 44 \hspace{0.8cm} 48 \hspace{0.8cm} 52 \hspace{0.8cm} 56 

\vspace{-2mm}
	
\hspace{1.3cm} 7160 \hspace{0.4cm}  4730 \hspace{0.5cm}  9320 \hspace{0.5cm}  9980 \hspace{0.5cm}  6950 \hspace{0.6cm} 6640
	\hspace{0.6cm}  7040 \hspace{0.5cm}  5130 \hspace{0.5cm} 7700
	\hspace{0.5cm}  7940 \hspace{0.5cm}  8000}
\vspace{1cm}



\vspace{1cm}

Slow neutron-capture process: \\	
\begin{center}
	\ce{C ->[He][n] O}
		
		%->[][e+,$\nu$_e]F ->[H][He]O ->[H][$\gamma$]F ->[H][$\gamma$]Ne ->[ ][e+,$\nu$_e]F}
\end{center}
\vspace{-2mm}
\hspace{7.7cm} {\scriptsize 13 \hspace{0.7cm} 16 \hfill \ce{^{13}C($\alpha$, n)^16O} \\
	
	\vspace{-2mm}

	\hspace{7.6cm} 6410 \hfill (Energy release in KeV)}

\vspace{1cm}

\begin{center}
	\ce{Ne ->[He][n] Mg}
	
	%->[][e+,$\nu$_e]F ->[H][He]O ->[H][$\gamma$]F ->[H][$\gamma$]Ne ->[ ][e+,$\nu$_e]F}
\end{center}
\vspace{-2mm}
\hspace{7.6cm} {\scriptsize 22 \hspace{0.9cm} 25 \hfill \ce{^{22}Ne($\alpha$, n)^25Mg} \\
	
	\vspace{-2mm}
	
	\hspace{7.6cm} 6410 \hfill (Energy release in KeV)}

\vspace{1cm}

rapid neutron-capture process: \\	
\begin{center}
	\ce{C ->[He][n] O}
	
	%->[][e+,$\nu$_e]F ->[H][He]O ->[H][$\gamma$]F ->[H][$\gamma$]Ne ->[ ][e+,$\nu$_e]F}
\end{center}
\vspace{-2mm}
\hspace{7.7cm} {\scriptsize 13 \hspace{0.7cm} 16 \hfill \ce{^{13}C($\alpha$, n)^16O} \\
	
	\vspace{-2mm}
	
	\hspace{7.6cm} 6410 \hfill (Energy release in KeV)}

\vspace{1cm}

\begin{center}
	\ce{Ne ->[He][n] Mg}
	
	%->[][e+,$\nu$_e]F ->[H][He]O ->[H][$\gamma$]F ->[H][$\gamma$]Ne ->[ ][e+,$\nu$_e]F}
\end{center}
\vspace{-2mm}
\hspace{7.6cm} {\scriptsize 22 \hspace{0.9cm} 25 \hfill \ce{^{22}Ne($\alpha$, n)^25Mg} \\
	
	\vspace{-2mm}
	
	\hspace{7.6cm} 6410 \hfill (Energy release in KeV)}

\vspace{1cm}


These are the some processes with which star generates an immense amount of energy to halt the gravitational collapse. If the star acquires a state of equilibrium in between gravitational collapse and thermal pressure, the state is named as hydro-static equilibrium. To express this state mathematically, gravitational pull and thermal pressure on the shell mass need to estimated as follows:

\begin{equation}
dm = 4 \pi r^2 \rho dr
\end{equation}

\begin{equation}
dP = - \frac{G M_r}{r^2} dr
\end{equation}

Similar to hydro-static equilibrium, energy transportation is also another important aspect which we need to consider. The energy generated by nuclear reaction near the core region will get transported outwards. If $\varepsilon$ is rate of energy generation per unit mass per unit time, the luminosity (energy flux per unit time) at radius r can be expressed as: 

\begin{equation}
dL_r = 4 \pi r^2 \rho \varepsilon dr
\end{equation}

The outward energy flux arises due to the temperature gradient.

\begin{equation}
dT = - \frac{G m \rho T	}{r² P} \ dr
\end{equation}

\subsection{Cepheid}

Cepheid of small period (P < 10 days) are tending towards F5-F8 spectral type at maximum brightness while appear to G or K type at minimum brightness. \cite{1996JRASC..90...82T} 


Cepheid are classified in the two main subclasses: Type I Cepheid (\textit{Classical Cepheid}) belongs to population I stars (high metallicity) with mass 4-20 $M_0$ and having period range 1 day to weeks with varing luminosity upto 100,000 $L_0$; and Type II Cepheid are population II star with longer periodicity upto 100 days. Type II Cepheid are also known as W Virginis stars. 

Classical Cepheid are yellow super gaints of spectral class F-6 to K-2 With the increasing age, pulsation of Cepheid decrease which makes them a good candidate for study star-formation history. \footnote{Bono, G. et al. Classical Cepheid pulsation models. X. The period-age relation. Astrophys. J. 621, 966—977 (2005)} 

After the discovery of period-luminosity relation \footnote{Leavitt, Henrietta S.; Pickering, Edward C. (1912)}, it become a general question that why does
Cepheids pulsate? To understand the pulsation of Cepheid, let us first recall why does a star
shine at the first place? A star shines because it burns hydrogen at its core and convert it
into heavier elements like helium. Based on thermodynamics of heat engine, Astronomer A. Eddington explained the mechanism of pulsation
of Cepheid by purposing a valve mechanism which is known as Kappa mechanism. It follows as when hydrogen burns out, helium becomes the next abundant element
on the surface of Red Giant star. Due to high surface temperature, Helium on the surface releases it both electrons and get ionized \ce{He^{2+}} . \ce{He^{2+}} is opaque and blocks the outcoming radiation, making the star appear dimmer. The
blocked radiation creates pressure which expands the star surface. Gained surface area increases
the luminosity of star meanwhile the expansion of surface also cools down the opaque He2+ . The
ionized Helium gets back to its ground state He which is a transparent gas. Consequently, all
the trapped radiation releases making the star appear even more luminous. As the radiation
pressure decreases, star shrinks down to original state and the process restarts again. To construct a model for the mechanism, we need to acknowledge the process of stellar evolution and its different stages. 

\subsection{Radial Pulsation and structure}

\begin{equation}
\dddot{\zeta} - 4 \pi r \dot{\zeta} \frac{\partial}{\partial m}[(3\Gamma_1 - 4)P] - \frac{1}{r^2}\frac{\partial}{\partial m}\left[16 \pi^2 \Gamma_1 P \rho r^6 \frac{\partial \dot\zeta}{\partial m}\right] + 4\pi r \frac{\partial}{\partial m} \left[\rho(\Gamma_3 - 1)\delta\left(\epsilon_{eff}-\frac{\partial L}{\partial m}\right)\right] = 0
\end{equation}
This is Linear non-adiabatic wave equation (LNAWE) \footnote{rosseman 1949}

\vspace{4cm}
